% 固定顺序检验；失败后停止，不再试探后续阈值。
\begin{tikzpicture}[
 x=1mm,y=1mm,font=\figurefont\small,
 process/.style={draw=BookBlue,fill=BookBlue!5,line width=0.85pt,
                 minimum width=51mm,minimum height=10mm,align=center},
 output/.style={draw=BookTeal,fill=BookTeal!5,line width=0.85pt,
                minimum width=34mm,minimum height=14mm,align=center},
 flow/.style={->,>=latex,draw=BookInk,line width=1.1pt,
              rounded corners=1.2mm}
]
 \path[use as bounding box] (0,0) rectangle (112,85);
 \node[process] (start) at (35,77) {固定模型、规则与阈值顺序};
 \node[process] (sets) at (35,58) {按当前阈值生成校准集合};
 \node[process] (check) at (35,39) {判断：$U_n(\lambda)\leq\alpha$？};
 \node[process] (save) at (35,17) {保存通过的规则};
 \node[output] (stop) at (94,39) {停止\\返回保存规则};
 \draw[flow] (start.south) -- (sets.north);
 \draw[flow] (sets.south) -- (check.north);
 \draw[flow] (check.south) -- node[right,font=\figurefont\footnotesize]
    {通过} (save.north);
 \draw[flow] (check.east) -- node[above,font=\figurefont\footnotesize]
    {失败} (stop.west);
 \draw[flow] (save.west) -- (3,17) -- (3,58) -- (sets.west);
 \node[text=BookMuted,font=\figurefont\footnotesize,rotate=90]
    at (0,38) {仍有阈值：继续};
 \draw[flow] (save.east) -- (94,17) -- (stop.south);
 \node[text=BookMuted,font=\figurefont\footnotesize] at (81,10)
    {序列用尽};
\end{tikzpicture}
